% !TEX root = ../main.tex
% ---------------------------------------------------------------------
%  P3 -- DMA
% ---------------------------------------------------------------------
\vlsec{cP3}{P3}{DMA (Direct Memory Access)}
\label{sec:dma}

\SW{Motivation} Ein Interrupt pro Byte kostet $\sim$100 Cycles
Kontextwechsel -- \emph{für ein Byte}. DMA: einmal konfigurieren, \emph{ein}
Interrupt für N Byte.

\begin{bsp}[4\,kB von einer Peripherie-FIFO]
byteweiser IRQ: $4096\times100=409\,600$ Cycles CPU-Zeit\\
DMA: $\sim$200 Setup $+$ 1 IRQ $\approx$ 300 Cycles \ra Faktor $\sim$1400.
\end{bsp}

\begin{formel}[Worin der Gewinn wirklich liegt]
Nicht in der Transferdauer -- die bestimmt die Peripherie -- sondern
darin, dass die CPU in dieser Zeit rechnen oder \emph{schlafen} kann.
Darum gehört DMA ins Fach „Power \emph{and} Performance“.
\end{formel}

\subsection{DMA als zweiter Busmaster}
\begin{ztab}{@{}lX@{}}
\R{BR}/\R{BG} & Bus Request/Grant: CPU schaltet ihre Treiber hochohmig \\
\R{RS}/\R{DS} & CPU adressiert und beschreibt die DMA-Register \\
\R{REQ}/\R{ACK} & Handshake mit dem I/O-Gerät \\
\R{IRQ} & Block fertig \\
\end{ztab}
\begin{formel}[Doppelrolle]
Beim \emph{Konfigurieren} ist der DMA \T{Slave}, beim \emph{Transfer}
\T{Master}. Genau das unterscheidet ihn von jeder anderen Peripherie.
(ARM DMAC: AXI-64-Master für Daten, APB-32-Slave für Register.)
\end{formel}

\SW{Arbitrierung}
\begin{ztab}{@{}lXl@{}}
Modus & Verhalten & CPU-Latenz \\ \midrule
Burst/Block & Bus bis Blockende halten & schlecht \\
Cycle Stealing & 1 Zyklus, freigeben, wieder & gut \\
Transparent & nur ungenutzte Zyklen & keine \\
\end{ztab}

\SW{Richtungen} mem\,\ra\,mem (voller Bustakt);
mem\,\ra\,Peripherie und Peripherie\,\ra\,mem (\emph{von der Peripherie
getaktet}, über full/empty-Flags oder req/ack);
externes Interface\,\ra\,mem mit Paketlänge $+$ IRQ (Ethernet).

\begin{formel}[Dasselbe Pacing-Konzept, drei Namen]
\T{Peripheral Request} (\R{DMAWFP}, PL330) $=$
\T{DREQ} (BCM2835) $=$ \T{TREQ} (RP2350).
\end{formel}

\subsection{Die zwei Fallstricke}
\label{falle:dma}
\begin{falle}[1 -- Cache-Kohärenz]
Write-Back heisst: das RAM ist veraltet, der externe DMA sieht den Cache
nicht.\\
\T{CPU\,\ra\,Peripherie:} Puffer noch dirty im Cache \ra DMA liest alte
Daten \ra vorher \T{clean}.\\
\T{Peripherie\,\ra\,CPU:} CPU trifft die alte Cache-Zeile \ra vorher
\T{invalidate}.\\
Alternativen: nicht cachebarer Puffer, oder HW-Kohärenz
(Zynq \T{ACP}-Port über die Snoop Control Unit).
\end{falle}

\begin{falle}[2 -- Physische Adressen]
Der DMA hat \emph{keine} MMU, keinen TLB. Ein Programm-Pointer ist
virtuell \ra stille Datenkorruption statt Segfault.\\
Und: ein virtuell zusammenhängender Puffer ist physisch fast nie
zusammenhängend (1\,MB $=$ 256 beliebig verteilte Frames).\\
\T{Daraus folgt die gesamte restliche Komplexität}: Control Blocks,
verkettete Listen und Befehlssätze existieren, damit ein Transfer aus
mehreren Stücken bestehen kann (\T{Scatter-Gather}). Alternative:
physisch zusammenhängend allozieren (CMA, VideoCore-Mailbox).
\end{falle}

\subsection{Zynq: Bandbreite}
\begin{formel}[Bandbreite]
\label{form:bw}
\[\text{Bandbreite}=\text{Takt}\times\text{Busbreite in Byte}\]
DMA läuft auf \R{CPU\_2x}. Bei CPU\_6x $=$ 800\,MHz (1:3) also 266\,MHz,
Central Interconnect 64\,Bit:
$266\,\text{MHz}\times8\,\text{B}=\T{2.13\,GB/s}$.
Slave Interconnect nur 32\,Bit \ra die Hälfte.
\end{formel}
Die Register hängen am APB mit \R{CPU\_1x} ($\sim$133\,MHz, nicht
gepipelinet) \ra \emph{ein} Registerwrite kostet Dutzende CPU-Cycles.
Darum lohnt DMA erst ab grossen Blöcken.
Clock-Domain-Übergänge (ASYNC, FIFO) kosten Latenz \ra \T{Burst} --
dasselbe Amortisationsargument wie die Cache Line in P1.

\subsection{ARM PL330: DMA mit eigenem Befehlssatz}
\begin{itemize}
\item AMBA-PrimeCell-IP von ARM; die Engine \emph{holt ihr Programm aus
      dem Speicher} über das AXI-Master-Interface
\item 8 Kanäle $=$ Threads, Round-Robin, gleiche Priorität
\item \T{eigener I-Cache}: 8 Lines à 4 Worte; während eines Line-Loads
      darf ein anderer Kanal aktiv sein
\item MFIFO $128\times64$\,Bit, von allen Kanälen geteilt --
      \emph{entkoppelt} Burstlänge und Breite von Quelle und Ziel
\item APB-Slave non-secure \R{0xF800\_4000} / secure \R{0xF800\_3000}
\end{itemize}

\begin{lstlisting}
DMAMOV CCR, SB4 SS64 DB4 DS64  # Src+Dst: 4 Beats a 64b
DMAMOV SAR, 0x1000             # Quelladresse
DMAMOV DAR, 0x4000             # Zieladresse
DMALP 16
  DMALD                        # 4 x 64b = 32 B laden
  DMAST                        # 32 B speichern
DMALPEND
DMAEND
\end{lstlisting}

\begin{rezept}[Rezept: \R{CCR} dekodieren]
\label{rez:ccr}
\R{S}/\R{D} $=$ Source/Destination, \R{B}$n$ $=$ Burst Length (Beats),
\R{S}$m$ $=$ Burst Size (Bit pro Beat).
\begin{enumerate}
\item Bytes pro \R{DMALD} $=n\times m/8$
\item Gesamtmenge $=$ Bytes pro Burst $\times$ Schleifenzahl
\item Prüfen, ob Quell- und Zielseite gleich viel bewegen
\end{enumerate}
Oben: $4\times64/8=32$\,B, $16\times32=\T{512\,B}$.
\end{rezept}

\begin{falle}[„Loop fits in one cache line (4 I-words)“]
Der Schleifenkörper hat genau 4 Instruktionen \ra einmal geholt,
16$\times$ ausgeführt. Ohne Source-Burst (\R{SB1}) braucht man
$4\times$\R{DMALD} $+$ \R{DMAST} $=$ 6 Instruktionen \ra passt nicht mehr,
die Engine fetcht in \emph{jedem} Durchlauf nach. Zweiter Verlust:
4 AXI-Adressphasen statt einer.\\
\T{Das ist Cache Trashing aus P1, eine Ebene tiefer.}
\end{falle}

\SW{Peripherie-Transfer} \R{DMAWFP} blockiert, bis die Peripherie ihren
Request zieht -- die anderen 7 Kanäle laufen weiter. Der Request trägt
einen \emph{Typ} (\R{DxTYPE}: 00 single, 01 burst); danach laufen Befehle
mit Suffix \R{B} nur bei Burst, mit \R{S} nur bei Single.
\T{Prädikation in der Zeit}, analog zu NEONs Masken im Raum.
\R{DMALPFE} $=$ Loop Forever, Abbruch über \R{DRLAST} der Hardware.
\R{DMASTP} $=$ Store \emph{and notify} (das ACK) \ra die letzte
Transaktion wird aus der Schleife gezogen (15$+$1, Loop Peeling).

\begin{ztab}{@{}lX@{}}
Daten & \R{DMALD/ST}, \R{DMALDP/STP}, \R{DMASTZ} (memset in HW) \\
Fluss & \R{DMALP/LPEND/LPFE}, \R{DMAGO/KILL/END/NOP} \\
Sync & \R{DMAWFP}; \R{DMAWFE}/\R{DMASEV} \emph{Kanal-zu-Kanal}
  \ra CPU-freie Pipelines; \R{DMAFLUSHP} \\
Barrieren & \R{DMARMB} / \R{DMAWMB} \\
Adressen & \R{DMAADDH} / \R{DMAADNH} \ra 2D und Stride \\
\end{ztab}

\begin{formel}[VALID/READY -- das universelle AMBA-Handshake]
Sender zieht \R{VALID}, Empfänger \R{READY}; die Übertragung findet im
Takt statt, in dem \emph{beide} high sind. Keine Seite darf auf die andere
warten, bevor sie selbst setzt (sonst Deadlock). Funktioniert bei
beliebigem Geschwindigkeitsverhältnis.
\end{formel}

\subsection{BCM2835 (Raspberry Pi): Control Blocks}
16 Kanäle, keine ISA, stattdessen eine \T{verkettete Liste} von CBs --
8 Worte à 32\,Bit $=$ 32\,B, 32-B-aligned, als \emph{ein} Burst geholt:
\begin{ztab}{@{}clcl@{}}
\# & Feld & \# & Feld \\ \midrule
0 & \R{TI} & 3 & \R{TXFR\_LEN} \\
1 & \R{SOURCE\_AD} & 4 & \R{STRIDE} (2D) \\
2 & \R{DEST\_AD} & 5 & \R{NEXTCONBK} \\
\multicolumn{4}{@{}l}{6--7 reserved, auf 0 setzen} \\
\end{ztab}
\SW{Start} CB-Adresse in \R{CONBLK\_AD} schreiben.
\SW{Stop} Transfer fertig \emph{und} \R{NEXTCONBK}$=0$ \ra ACTIVE-Bit wird
gelöscht.
\SW{Warum 8 Worte} 32\,B ist eine natürliche Burstgrösse und eine
Zweierpotenz \ra CBs sind array-indizierbar und ragen nie über eine Grenze.
\SW{\R{NEXTCONBK}} erlaubt Scatter-Gather und Ringpuffer (letzter CB zeigt
auf den ersten) komplett ohne CPU.

\begin{rezept}[Rezept: \R{TI}-Register setzen]
\label{rez:ti}
Zeigt die Adresse auf \emph{Speicher}, muss sie hochzählen; zeigt sie auf
ein \emph{Peripherieregister} (FIFO), darf sie sich \textbf{nicht} ändern.
\begin{tabularx}{\linewidth}{@{}Xccc@{}}
Transfer & \R{SRC\_INC} & \R{DEST\_INC} & \R{DREQ} \\
mem \ra mem & 1 & 1 & -- \\
mem \ra Periph. & 1 & \textbf{0} & DEST \\
Periph. \ra mem & \textbf{0} & 1 & SRC \\
\end{tabularx}\\[2pt]
Vergessen $=$ der DMA überschreibt den halben Peripherie-Adressraum.
\end{rezept}

\begin{ztab}{@{}lX@{}}
\R{PERMAP} 20:16 & welches Ready-Signal (1--31) \\
\R{WAITS} 25:21 & Wartezyklen einfügen $=$ Bandbreite drosseln \\
\R{BURST\_LENGTH} 15:12 & Burstlänge \\
\R{SRC/DEST\_WIDTH} & 32 oder 128\,Bit \\
\R{WAIT\_RESP} 3 & auf AXI Write Response warten \\
\R{TDMODE} 1 & 2D-Modus, nutzt \R{STRIDE} \ra Rechteck aus Framebuffer \\
\R{INTEN} 0 & Interrupt freigeben \\
\R{SRC\_IGNORE} 11 & nicht lesen \ra schnelles Nullen (memset) \\
\R{DEST\_IGNORE} 7 & nicht schreiben \ra \emph{Prefetch in den L2} \\
\end{ztab}

\begin{formel}[Panic-Signal $=$ QoS in Hardware]
Droht einer Peripherie-FIFO Under- oder Overflow, erhöht der DMA
\emph{dynamisch} sein AXI-Prioritätslevel. Normalbetrieb niedrige
Priorität (CPU nicht stören), Notfall Vorfahrt. Damit muss man sich nicht
zwischen Durchsatz und CPU-Latenz entscheiden.
\end{formel}

\subsection{RP2350: vier Register pro Kanal}
16 Kanäle, kein CB, keine ISA: \R{READ\_ADDR}, \R{WRITE\_ADDR},
\R{TRANS\_COUNT}, \R{CTRL} -- alle vier zählen während der Ausführung mit
(Fortschritt jederzeit lesbar).
\SW{Dual bus ports} je ein Lese- \emph{und} ein Schreibzugriff pro Takt
\ra 1 Transfer/Takt im Dauerbetrieb.

\begin{formel}[Der RP2350-Trick]
Die DMA-Konfigurationsregister sind normale memory-mapped Register --
also kann \emph{ein DMA-Kanal einen anderen programmieren}. Damit baut man
Scatter-Gather aus einem Design ohne jede Listenunterstützung nach.
\end{formel}

\begin{ztab}{@{}lX@{}}
\R{EN}, \R{HIGH\_PRIORITY} & aktiv, mehr Bandbreite \\
\R{DATA\_SIZE} 3:2 & 8/16/32\,Bit \\
\R{INCR\_READ/WRITE} ($+$\R{\_REV}) & hochzählen / rückwärts \\
\R{RING\_SEL}, \R{RING\_SIZE} & „freeze all but $n$ lower bits“
  $=$ \T{Ringpuffer in HW}, Puffer auf $2^n$ ausgerichtet \\
\R{CHAIN\_TO} 16:13 & anderen Kanal auslösen \\
\R{TREQ\_SEL} 22:17 & DREQ-Nummer oder Timer \\
\R{BSWAP} 24 & Endianness gratis \\
\R{SNIFF\_EN} 25 & \T{CRC während des Transfers} \\
\R{BUSY} 26, Fehlerflags 31:29 & Status \\
\end{ztab}
\SW{Weiter} \R{TRANS\_COUNT.MODE} (\R{0x0} stoppen, \R{0x1} neu laden
$+$ IRQ, \R{0xF} endlos); \R{INTR/INTE/INTF/INTS};
\R{TIMER0-3} (Taktteiler \ra Transfer mit exakt fester Rate ganz ohne
Peripherie, z.\,B.\ 44.1\,kHz an einen DAC);
\R{SNIFF\_CTRL} (CRC-32 IEEE\,802.3 / CRC-16-CCITT / 32-Bit-Summe);
\R{MULTI\_CHAN\_TRIGGER}; \R{CHAN\_ABORT}.

\subsection{Drei Philosophien im Vergleich}
\begin{ztab}{@{}lXXX@{}}
 & PL330 & BCM2835 & RP2350 \\ \midrule
Beschreibung & Programm & CB-Liste & 4 Register \\
Kanäle & 8 & 16 & 16 \\
Verkettung & Programm & \R{NEXTCONBK} & \R{CHAIN\_TO} \\
Schleifen & ja & nein & nein \\
Komplexität & hoch & mittel & niedrig \\
\end{ztab}
Drei Punkte auf dem Spektrum Zustandsmaschine\,$\leftrightarrow$\,Prozessor.
\emph{Alle drei lösen dasselbe Problem} -- wer Falle 2 verstanden hat,
kann sich jede der drei Lösungen herleiten.
